% Generated draft: requires mathematical and source review.
\documentclass[11pt]{article}
\usepackage[margin=25mm]{geometry}
\usepackage{fontspec}
\setmainfont{Latin Modern Roman}
\usepackage{amsmath,amssymb,mathtools,unicode-math}
\setmathfont{Latin Modern Math}
\usepackage{enumitem,needspace,xurl,hyperref}
\setlength{\parindent}{0pt}
\setlength{\parskip}{5pt}
\setlength{\emergencystretch}{3em}
\newcommand{\N}{\mathbb{N}}
\newcommand{\Z}{\mathbb{Z}}
\newcommand{\Q}{\mathbb{Q}}
\newcommand{\R}{\mathbb{R}}
\newcommand{\C}{\mathbb{C}}
\newcommand{\F}{\mathbb{F}}
\newcommand{\A}{\mathbb{A}}
\newcommand{\PP}{\mathbb{P}}
\newcommand{\E}{\mathbb{E}}
\newcommand{\eps}{\varepsilon}
\newcommand{\dd}{\,\mathrm{d}}
\begin{document}
\section*{UnsolvedMath ID: 10400187}
\textbf{Source problem code:} catalogue entry
\subsection{Definitions and mathematical statement}
Let \(\mathcal H_{\Z}\) be the class of integral homology \(3\)-spheres, conventionally closed oriented \(M\) with \(H_i(M;\Z)\cong H_i(S^3;\Z)\) for every \(i\). A finite-type invariant is understood relative to a specified surgery filtration \(F_\bullet\) on the linear span of this class: type at most \(d\) means that its linear extension vanishes on \(F_{d+1}\). The source extract does not define that filtration, coefficient system or degree convention.

For a manifold \(M\), the ordered configuration space of \(r\) points is
\[
 \operatorname{Conf}_r(M)=
 \{(x_1,\ldots,x_r)\in M^r:x_i\ne x_j\text{ whenever }i\ne j\}.
\]
Configuration-space integral constructions use appropriate compactifications, differential forms and boundary/correction data; none of those choices or formulas is given here.

Problem 10.12 asks for a constructive combinatorial presentation of \emph{each} finite-type invariant of integral homology \(3\)-spheres, obtained by localizing configuration-space integrals. It singles out the Casson invariant as a particular target. The requested method must relate the localized integral construction to the resulting combinatorial formula, not merely provide an unrelated existing formula for the same invariant. The source gives no localization mechanism, permitted combinatorial input, uniform algorithm, or normalization for the Casson invariant. These must be supplied to make the construction and its invariance verifiable; no choice of propagator, compactification or surgery degree is imposed here.
\subsection{Short English statement}
Produce constructive combinatorial formulas for every finite-type invariant of integral homology three-spheres by localizing their configuration-space integrals, especially for the Casson invariant. The record does not specify the integral model, localization procedure or finite-type conventions.
\subsection{Sources}
\begin{itemize}
\item Local UnsolvedMath registry, permanent ID \(10400187\), exact source code \texttt{catalogue entry}, Problem 10.12. No individual proposer or proposed year is supplied; the title is truncated.
\item Ohtsuki, \emph{Problems on invariants of knots and 3-manifolds} (2002), Problem 10.12, PDF page 150: \url{https://msp.org/gtm/2002/04/gtm-2002-04-024s.pdf}. Bold headings distinguish 231 numbered items from references in surrounding exposition.
\item Extraction pdf-font-structure, accessed 2026-07-29. Source evidence NEEDS\_REVIEW presents an open problem and requires a current primary-literature audit. Supplied classification OPEN-TRIAGE identifies the configuration-integral localization program, with no further primary references or resolution claim.
\item Dataset attribution: Ulam AI / UnsolvedMath Contributors, CC BY 4.0. Source rights NEEDS\_REVIEW: publicly accessible author/publisher copy, redistribution terms requiring review.
\end{itemize}
\end{document}
